\chapter{Two insertions and one extension}
\label{cyopen:chapter}

An extension distinguishes the subobject from the quotient.
This order can remain visible after passing to cohomology.
For the quiver with one vertex and one weighted loop, reversing
the two inputs changes their product.
The same geometry has a nonsplit extension at coincident support.
Its invariant lines form a nonreduced scheme of length two.

\section{The first ordered product}

Geometric objects are over $\mathbb C$, and cohomology has rational
coefficients. Put $R=\mathbb Q[\hbar]$ and
\[
 \mathcal S_n=R[x_1,\ldots,x_n]^{\mathfrak S_n},\qquad
 \mathcal S=\bigoplus_{n\geq0}\mathcal S_n.
\]
The symmetric group permutes the variables. The summand index $n$
is called the \emph{charge}. Write $\mathbf1_n$ for the constant
polynomial in $\mathcal S_n$, including $\mathcal S_0=R$.

Give every $x_i$ and $\hbar$ cohomological degree two.
All classes are even. The coefficient ring is polynomial, with no
completion.

For $f,g\in\mathcal S_1$, define
\begin{equation}\label{cyopen:eq:binary}
 \begin{split}
 (f\star g)(x_1,x_2)
 ={}&f(x_1)g(x_2)\frac{x_2-x_1-\hbar}{x_2-x_1}\\
 &+f(x_2)g(x_1)\frac{x_1-x_2-\hbar}{x_1-x_2}.
 \end{split}
\end{equation}
The two summands correspond to the two choices of a first root.
Their residues at $x_1=x_2$ cancel, so the result is a symmetric
polynomial.

Set $a_r=x_1^r\in\mathcal S_1$ for $r\geq0$.
Write $[f,g]=f\star g-g\star f$.
Direct substitution gives
\begin{equation}\label{cyopen:eq:first-products}
 \begin{aligned}
 a_0\star a_1&=x_1+x_2-\hbar,&
 a_1\star a_0&=x_1+x_2+\hbar,\\
 a_0\star a_0&=2\mathbf1_2,&
 [a_1,a_0]&=2\hbar\mathbf1_2.
 \end{aligned}
\end{equation}
In particular, $[a_1,a_0]=\hbar(a_0\star a_0)$.
The constant $\mathbf1_2$ has charge two. The algebra unit will be
$\mathbf1_0$, and $a_0=\mathbf1_1$ is a third element.

The two-root formula extends by assigning roots to the two inputs.
For a subset $I$, write $x_I$ for its list of roots.
For $f\in\mathcal S_n$ and $g\in\mathcal S_m$, set
\begin{equation}\label{cyopen:eq:shuffle}
 (f\star g)(x)=
 \sum_{\substack{I\subset\{1,\ldots,n+m\}\\|I|=n}}
 f(x_I)g(x_{I^c})
 \prod_{i\in I,\,j\in I^c}
 \frac{x_j-x_i-\hbar}{x_j-x_i}.
\end{equation}
Symmetry makes the order within $I$ and $I^c$ irrelevant.
An empty product equals one. Each subset occurs once.

\begin{proposition}[The weighted-loop algebra]
\label{cyopen:prop:algebra}
Formula~\eqref{cyopen:eq:shuffle} defines an associative unital
$R$-algebra structure on $\mathcal S$. Multiplication adds charges
and has cohomological degree zero. Its unit is $\mathbf1_0$.
\end{proposition}
\begin{proof}
The only possible poles are simple poles on root diagonals.
Fix $x_i=x_j$ and pair each summand that places these roots in
opposite blocks with the summand that exchanges them.
On this diagonal the input polynomials and all other pair factors
agree, while the two residues have opposite signs.
Thus every residue vanishes.

The common denominator is a product
of distinct diagonal equations in the unique factorization domain
$R[x_1,\ldots,x_{n+m}]$.
Cancellation at each one proves polynomiality.
Permuting the roots permutes the summands.

For three inputs, either bracketing sums over the same ordered
partition of the roots into three blocks. Every pair of roots in
distinct blocks contributes its factor exactly once, in the order
of those blocks. The two rational expressions agree term by term
after this common indexing, and hence agree as polynomials.
With an empty input there is one subset and no pair factor,
which proves the unit identities. Each pair factor has degree zero.
\end{proof}

At $\hbar=0$ the pair factors are one. The commutator in
\eqref{cyopen:eq:first-products} therefore measures a deformation
of the unweighted product. Its geometric origin is the difference
between the loop weight and the tangent weight of a flag.

\section{The Euler factor of an invariant subspace}

A representation of the one-loop quiver is a vector space with
an endomorphism. In dimension $n$ its moduli stack is
\[
 \mathfrak M_n=[\operatorname{End}(\mathbb C^n)/GL_n].
\]
Let $T=\mathbb C^\times$ act on the endomorphism by the character
$\chi_{-1}(s)=s^{-1}$, and let $\hbar$ be the first Chern class
of the character $s\mapsto s$. Equivariant contraction of the
endomorphism space to zero identifies
\[
 H^\bullet_T(\mathfrak M_n;\mathbb Q)
 =H^\bullet(BGL_n\times BT;\mathbb Q)=\mathcal S_n.
\]
The variables $x_i$ are the Chern roots of the universal vector
bundle. They have cohomological degree two, independently of the
eigenvalues of an individual endomorphism.

Let $\mathfrak E_{n,m}$ classify an endomorphism of an
$(n+m)$-dimensional space together with an invariant
$n$-dimensional subspace. It gives the correspondence
\begin{equation}\label{cyopen:eq:extension-correspondence}
 \mathfrak M_n\times\mathfrak M_m
 \xleftarrow{\ q\ }\mathfrak E_{n,m}
 \xrightarrow{\ p\ }\mathfrak M_{n+m}.
\end{equation}
The map $q$ takes the restricted endomorphism and its quotient.
The map $p$ forgets the subspace. It is proper, since invariant
subspaces form a closed subscheme of a Grassmannian.

On $\operatorname{Gr}(n,n+m)$ write $S$ for the tautological
subbundle and $Q$ for the quotient bundle. Invariance sets to
zero the endomorphism block from $S$ to $Q$.
Its normal bundle is
$\operatorname{Hom}(S,Q)\otimes\chi_{-1}$.
The Grassmannian tangent bundle is $\operatorname{Hom}(S,Q)$.

Write $e$ for the top Chern class and $f(S),g(Q)$ for evaluation
on the respective Chern roots.
The Hall operation in ordinary equivariant cohomology is
complex-oriented integration of
\begin{equation}\label{cyopen:eq:euler-integration}
 f(S)g(Q)\,
 e\bigl(\operatorname{Hom}(S,Q)\otimes\chi_{-1}\bigr)
\end{equation}
over the Grassmannian.
The normal Gysin map adds degree $2nm$ and Grassmannian
integration subtracts $2nm$.

At the coordinate subspace indexed by $I$, these two bundles have
roots $x_j-x_i-\hbar$ and $x_j-x_i$, respectively, for
$i\in I$, $j\in I^c$. Equivariant integration divides by the
tangent Euler class and sums over these subspaces.
It gives exactly \eqref{cyopen:eq:shuffle}.
This is the ordinary quiver Hall construction of
\cite{cyopen:KS}*{\S\S2.4--2.5}, with the loop character retained.
Proposition~\ref{cyopen:prop:algebra} proves directly that its
formula is polynomial over $\mathbb Q[\hbar]$.

For $n=m=1$, $S=\mathcal O(-1)$ and $Q=\mathcal O(1)$
on $\mathbb P^1$. Consequently
\[
 \int_{\mathbb P^1}c_1(Q)=1,\qquad
 \int_{\mathbb P^1}c_1(S)=-1,\qquad
 \int_{\mathbb P^1}c_1\operatorname{Hom}(S,Q)=2.
\]
These equations fix the signs and the charge-two coefficient in
\eqref{cyopen:eq:first-products}. Geometrically, the coefficient
integrates invariant-subspace data. Those data change when support
points coincide.

\section{A nonsplit extension at a collision}

To see the support, regard an endomorphism as the action of $t$
on a finite-dimensional $A=\mathbb C[t]$-module.
For $\lambda\in\mathbb C$, put $L_\lambda=A/(t-\lambda)$.
An extension
\begin{equation}\label{cyopen:eq:extension}
 0\longrightarrow L_\lambda\longrightarrow E
 \longrightarrow L_\mu\longrightarrow0
\end{equation}
has, after choosing a lift of the quotient basis, the matrix
\[
 T_E=\begin{pmatrix}\lambda&c\\0&\mu\end{pmatrix}.
\]
Replacing that lift by itself plus $b$ times the submodule basis
changes $c$ to $c+b(\lambda-\mu)$.
If $\lambda\ne\mu$, taking $b=-c/(\lambda-\mu)$ splits the
extension. If $\lambda=\mu$, the parameter $c$ is unchanged.

The free resolution
$0\to A\xrightarrow{t-\mu}A\to L_\mu\to0$
computes the full derived Hom by the two-term complex
\begin{equation}\label{cyopen:eq:hom-complex}
 \operatorname{RHom}_A(L_\mu,L_\lambda)\simeq
 [\,\mathbb C\xrightarrow{\lambda-\mu}\mathbb C\,],
 \qquad\text{in degrees }0,1.
\end{equation}
The degree-one basis is chosen for the displayed differential.
Both Hom and $\operatorname{Ext}^1$ vanish at distinct supports.
At coincident support they are each one-dimensional.
In particular, the splitting off the diagonal is unique as a
splitting of \eqref{cyopen:eq:extension}.

At $\lambda=\mu=0$ and $c=1$ the middle term is $A/(t^2)$.
Its exact sequence is
\begin{equation}\label{cyopen:eq:jordan-extension}
 0\longrightarrow A/(t)
 \xrightarrow{\,1\mapsto t\,}A/(t^2)
 \longrightarrow A/(t)\longrightarrow0.
\end{equation}
It is nonsplit because multiplication by $t$ is nonzero on its
middle term and zero on $A/(t)\oplus A/(t)$.
The characteristic polynomial of $T_E$ is nevertheless
$(s-\lambda)(s-\mu)$ for every $c$.
It records the two support points with multiplicity and forgets
the extension class.

The flag correspondence retains still more than the invariant
lines as a set. In the affine chart of $\mathbb P^1$ whose line
is spanned by $(1,w)$, invariance under $T_E$ is the equation
\begin{equation}\label{cyopen:eq:stable-line}
 (\mu-\lambda)w-cw^2=0.
\end{equation}
Indeed, the second coordinate of $T_E(1,w)$ must equal $w$
times the first. This equation holds over every commutative
$\mathbb C$-algebra, so it describes the scheme, including its
infinitesimal points.

\begin{proposition}[The invariant-line scheme of a Jordan block]
\label{cyopen:prop:double-line}
For $T_E=\left(\begin{smallmatrix}0&1\\0&0\end{smallmatrix}\right)$,
the scheme of invariant lines is
$\operatorname{Spec}\mathbb C[w]/(w^2)$.
It has length two and one underlying point.
\end{proposition}
\begin{proof}
The only invariant line over $\mathbb C$ is the first coordinate
line, so the entire invariant-line scheme lies in the displayed
affine chart. Equation~\eqref{cyopen:eq:stable-line} becomes $w^2=0$.
The quotient has basis $1,w$ over $\mathbb C$.
\end{proof}

Thus the two in \eqref{cyopen:eq:first-products} is also the
length of this regular zero scheme of the invariant-line section.
Replacing the scheme by its reduced point loses that length.
For the scalar endomorphism the section instead vanishes identically,
and the invariant-line scheme is all of $\mathbb P^1$.
The Euler-class operation remains the specified pushforward.

\section{Relations and ordered words}

The first commutator belongs to a complete presentation of the
weighted-loop algebra. Let $\mathcal P_\hbar$ be the unital
associative $R$-algebra generated by $e_r$, $r\geq0$, subject to
\begin{equation}\label{cyopen:eq:relations}
 [e_{r+1},e_s]-[e_r,e_{s+1}]
 =\hbar(e_re_s+e_se_r),\qquad r,s\geq0.
\end{equation}
Each $e_r$ has charge one and cohomological degree $2r$.

\begin{theorem}[An integral ordered basis]
\label{cyopen:thm:presentation}
The map $e_r\mapsto a_r$ is an isomorphism
$\mathcal P_\hbar\simeq\mathcal S$ of charge-graded $R$-algebras.
The words $e_{r_1}\cdots e_{r_n}$ with
$0\leq r_1\leq\cdots\leq r_n$, including the empty word,
form an $R$-basis.
\end{theorem}
\begin{proof}
For two roots $x,y$, formula~\eqref{cyopen:eq:binary} gives
\[
 [a_r,a_s]=2\hbar\frac{y^rx^s-x^ry^s}{y-x}.
\]
Subtracting this formula at $(r,s+1)$ from the one at $(r+1,s)$
gives $2\hbar(x^ry^s+x^sy^r)$.
This also equals $\hbar(a_r\star a_s+a_s\star a_r)$,
so the assignments respect \eqref{cyopen:eq:relations}.

These relations order the words. For $r>s$, write
$C_{r,s}=[e_r,e_s]$. Then
\[
 C_{r,s}=C_{r-1,s+1}
       +\hbar(e_{r-1}e_s+e_se_{r-1}).
\]
The index gap decreases until
$C_{s,s}=0$ or $C_{s+1,s}=\hbar e_s^2$.
Every correction has smaller sum of generator indices.
Induction on that sum, and then on the number of inversions in
the word, expresses every word in ordered form.

Filter $\mathcal S_n$ by total root degree, treating $\hbar$
as a coefficient. The leading part of an ordered product is
\[
 \sum_{\sigma\in\mathfrak S_n}
 x_1^{r_{\sigma(1)}}\cdots x_n^{r_{\sigma(n)}}.
\]
It is the monomial symmetric polynomial with these exponents
multiplied by the factorials of their multiplicities.
Every such factorial is invertible in $R$.
Distinct ordered words have distinct leading basis elements,
which proves independence over $R$.

Subtracting the leading term of a symmetric polynomial reduces
its root degree and proves spanning. The ordered spanning words
in $\mathcal P_\hbar$ therefore map bijectively to a basis of
$\mathcal S$.
\end{proof}

The proof works at $\hbar=0$ and over the full ring $\mathbb Q[\hbar]$.
It gives an algebraic deformation with its generators and relations.
A field realization requires operations with a different target.

\section{A coefficient transpose and the field comparison}

Let $m_{1,1}:\mathcal S_1\otimes_R\mathcal S_1\to\mathcal S_2$
be the binary product~\eqref{cyopen:eq:binary}.
Let $\epsilon_r:\mathcal S_1\to R$ extract the coefficient of
$x_1^r$, and let $\epsilon^{(2)}:\mathcal S_2\to R$ extract the
constant term in both roots. Transposing the binary Hall product
on this particular functional gives the finite tensor
\begin{equation}\label{cyopen:eq:dual-binary}
 m_{1,1}^*\epsilon^{(2)}
 =2\epsilon_0\otimes\epsilon_0
  +\hbar(\epsilon_1\otimes\epsilon_0-
                   \epsilon_0\otimes\epsilon_1).
\end{equation}

To verify the formula, evaluate on
$a_r\otimes a_s$. The polynomial part of the product contributes
a constant only for $r=s=0$.
Its part proportional to $\hbar$ contributes only when $r+s=1$.
Equation~\eqref{cyopen:eq:first-products} gives the three coefficients.

For an $R$-module $V$, put $V((z))=V[[z]][z^{-1}]$.
A binary field operation has type
\begin{equation}\label{cyopen:eq:field-type}
 Y:V\otimes_R V\longrightarrow V((z)).
\end{equation}
Thus $Y(v,z)w$ has a finite principal part for each pair $v,w$.
This target differs from the tensor of coefficient functionals
in \eqref{cyopen:eq:dual-binary}.
The formal variable $z$ describes an insertion difference.
The Chern roots in \eqref{cyopen:eq:binary} describe characteristic
classes on a moduli stack.

\begin{problem}[Realizing extension operations as fields]
\label{cyopen:prob:realization}
For specified Calabi--Yau moduli and extension correspondences,
construct their oriented coefficient complexes on finite support
configurations and their maps across collision strata.
Construct a state module, a vacuum, translation, and operations of type
\eqref{cyopen:eq:field-type}, together with a comparison
of their coefficients with the geometric operations.
The comparison must preserve the empty-object unit, the chosen
local vacuum, nested flag composition, and the attaching maps
at collisions. Its domain and its tensor completions are part of the construction.
\end{problem}

The one-loop model supplies the Hall multiplication and
\eqref{cyopen:eq:dual-binary}. Its coincident extension shows
what a comparison based only on support would omit.
For Calabi--Yau moduli, the coefficient construction additionally
uses the cyclic duality class and compatible orientation data.
These enter before the extension correspondences can define the
required operations on complexes.
